\section{Keep the Gram matrix and integrate directly}\label{sec:hua}

\subsection{The main source-shortening recommendation}

For the first source port, this is the stronger recommendation. The existing proof
constructs a negative frame, normalizes its Gram matrix to \(-\Id\) using a lower
triangular matrix \(L\), follows both \(\det L\) and its final diagonal entry,
and subsequently cancels those factors in a Hua integral. The final inequality
depends only on two invariants of the unnormalized frame. Retaining them directly
removes the purpose of the triangular normalization.

The central integral has a fixed final column. For an \(i+1\) by \(i+1\) positive
Hermitian matrix \(H\), a fixed \(v\in\C^k\), and a variable \(U\in\C^{k\times i}\),
write \([U,v]\) for the matrix formed by appending \(v\).

\begin{theorem}[Bordered matrix integral]\label{thm:hua}
Let \(k\geq1\), \(i\geq0\), and \(H=H^*>0\), of size \(i+1\).
Write \(h=H_{i+1,i+1}>0\), and use ordinary product Lebesgue measure
\(\prod_{a,b}d\Re U_{ab}\,d\Im U_{ab}\). Then
\begin{equation}\label{eq:hua}
 \int_{\C^{k\times i}}
   \det\bigl(H+[U,v]^*[U,v]\bigr)^{-(k+i)}\,dU
 =C_{k,i}\,(\det H)^{-i}\bigl(h+\norm v^2\bigr)^{-k},
\end{equation}
where
\begin{equation}\label{eq:hua-constant}
 C_{k,i}=\pi^{ki}\prod_{\ell=1}^{i}\frac{\Gamma(\ell)}{\Gamma(k+\ell)}
        =\pi^{ki}\prod_{\ell=1}^{i}\frac{(\ell-1)!}{(k+\ell-1)!}.
\end{equation}
For \(i=0\), the integral is over a one-point space, both products are empty,
and the same identity holds.
\end{theorem}

This is a direct specialization of the classical matrix beta-integral calculus.
The contribution to this project is the invariant formulation and its exact adapter
to the current majorant, not a claim that the general Hua-integral family is new.

\subsection{Proof with every determinant factor retained}

Partition \(H\) and define
\begin{equation}\label{eq:schur-data}
 H=\begin{pmatrix}H_0&b\\b^*&h\end{pmatrix},\quad
 S=H_0-\frac{bb^*}{h},\quad
 a=h+\norm v^2,\quad
 G=\Id-\frac{vv^*}{a},\quad
 U_c=\frac{vb^*}{h}.
\end{equation}
Here \(S>0\), \(G>0\), and
\begin{equation}\label{eq:schur-dets}
 \det S=\frac{\det H}{h},\qquad
 \det G=\frac{h}{a},\qquad
 G^{-1}=\Id+\frac{vv^*}{h}.
\end{equation}
The statement about \(G\) follows either from its eigenvalues or from the displayed
rank-one inverse.

Let \(W=U-U_c\). A unit-determinant block congruence, followed by a Schur complement,
gives
\begin{equation}\label{eq:complete-square}
 \det\bigl(H+[U,v]^*[U,v]\bigr)
       =a\det\bigl(S+W^*GW\bigr).
\end{equation}
For clarity, the congruence matrix is
\[
 C=\begin{pmatrix}\Id&0\\-b^*/h&1\end{pmatrix}.
\]
It satisfies \(C^*HC=\operatorname{diag}(S,h)\) and
\([U,v]C=[W,v]\). The lower-right entry of the resulting Gram matrix is
\(a\), so taking its Schur complement proves \eqref{eq:complete-square}.

Now put \(W=G^{-1/2}ZS^{1/2}\). Since the variables are complex matrices,
the real Jacobian is
\[
 dU=(\det G)^{-i}(\det S)^k\,dZ.
\]
Indeed, a complex linear map has real determinant equal to the squared modulus
of its complex determinant; left multiplication acts on \(i\) columns and
right multiplication on \(k\) rows.
Thus the left side of \eqref{eq:hua} becomes
\begin{align*}
 &a^{-(k+i)}(\det S)^{-(k+i)}
    (\det G)^{-i}(\det S)^k
    \int_{\C^{k\times i}}\det(\Id+Z^*Z)^{-(k+i)}\,dZ\\
 &\qquad=(\det H)^{-i}a^{-k}
    \int_{\C^{k\times i}}\det(\Id+Z^*Z)^{-(k+i)}\,dZ.
\end{align*}
The powers of \(h\) cancel exactly. This cancellation is the reason for keeping the
last column visible rather than absorbing it into a normalized basis.

To evaluate the remaining integral, let
\[
 J_{k,i}(p)=\int_{\C^{k\times i}}\det(\Id+Z^*Z)^{-p}\,dZ.
\]
Integrate the last column \(z\) after fixing the others \(X\). The determinant lemma
and a positive Hermitian change of variables give
\[
 \det(\Id+[X,z]^*[X,z])
 =\det(\Id+X^*X)
       \bigl(1+z^*(\Id+XX^*)^{-1}z\bigr).
\]
The scalar integral over \(\C^k\), computed in polar coordinates, is
\[
 \int_{\C^k}(1+\norm z^2)^{-p}\,dz
       =\pi^k\frac{\Gamma(p-k)}{\Gamma(p)}
       \qquad(p>k).
\]
Therefore, for \(p>k+i-1\),
\[
 J_{k,i}(p)=\pi^k\frac{\Gamma(p-k)}{\Gamma(p)}
              J_{k,i-1}(p-1),\qquad J_{k,0}(p)=1.
\]
Taking \(p=k+i\) yields \eqref{eq:hua-constant} and proves the theorem.
All integrands are nonnegative, so the successive integrations are justified
by Tonelli; the displayed recurrence also proves finiteness.

\begin{remark}[A useful normalization check]
For \(k=i=1\), take
\[
 H=\begin{pmatrix}3&1+\mathrm i\\1-\mathrm i&2\end{pmatrix},
 \qquad v=1+2\mathrm i.
\]
Then \(\det H=4\), \(h+\abs v^2=7\), and direct expansion gives
\[
 \det(H+[u,v]^*[u,v])
   =14+2\left|u-\frac{3+\mathrm i}{2}\right|^2.
\]
The integral of its inverse square over \(\C\) is \(\pi/28\),
exactly as \eqref{eq:hua} predicts. This example checks the conjugation in \(U_c\),
the complex Jacobian, and the exponent of the final-column factor.
\end{remark}

\subsection{The shorter Lean route uses an existing integral bound}

The exact integral is useful mathematically, but its full formalization is not
needed to shorten the source. The already available theorem
\src{Lib/General/Analysis/HolomorphicIntegrals}{lintegral\_norm\_det\_add\_inv\_pow\_le}
states, for \(R\in\C^{i\times i}\) invertible, \(B_0>0\) of size \(k\),
\(k\geq7\), and \(C_0\geq0\), that
\begin{equation}\label{eq:existing-hua}
 \int_{\C^{i\times k}}
 C_0\abs{\det(B_0+(RA+c_0)^*(RA+c_0))}^{-(k+i)}\,dA
 \leq C_0\abs{\det R}^{-2k}(\det B_0)^{-k}.
\end{equation}
It already includes the needed measure normalization and affine change of variables.

Choose any \(R\) with \(R^*R=S^{-1}\). Sylvester's determinant identity transforms
\eqref{eq:complete-square} into
\begin{equation}\label{eq:old-hua-adapter}
 \det(H+[U,v]^*[U,v])
  =\det H\,
     \det\bigl(G^{-1}+(RW^*)^*(RW^*)\bigr).
\end{equation}
The source variable is \(A\in\C^{i\times k}\), and \(U=U_{\rm base}+A^*\).
Hence
\[
 RW^*=RA+R(U_{\rm base}-U_c)^*.
\]
Apply \eqref{eq:existing-hua} with
\[
 B_0=G^{-1},\qquad c_0=R(U_{\rm base}-U_c)^*,\qquad
 C_0=(\det H)^{-(k+i)}.
\]
Because \(\abs{\det R}^2=(\det S)^{-1}\) and
\(\det(G^{-1})=a/h\), the bound becomes
\[
 (\det H)^{-(k+i)}(\det S)^k(h/a)^k
       =(\det H)^{-i}a^{-k}.
\]
Thus the source port needs no new gamma-function evaluation, no matrix-square-root
change-of-measure theorem, and no separate conjugate-transpose measure transport.
The only factorization of a positive matrix that remains is an unconstrained square
root of the Schur complement used at the integral interface; no triangular flag
or last-diagonal-entry bookkeeping is attached to it.

\subsection{The unnormalized frame and the volume ratio}

Here is the adapter to the geometry of \code{OneRealPlace}.
Let \(F\) be the full frame, \(J\) its ambient Hermitian involution, and
\(P_y=\code{frameLeftInverse}\,y\). Let \(E\) include the \(i+1\) subframe
coordinates, including the final anisotropic vector. Put
\begin{equation}\label{eq:raw-data}
 \begin{split}
 Z&=F^*J\gamma F,\qquad B=E^*ZE,\qquad D=ZE B^{-1},\\
 \omega_0&=JP_y^*D,\qquad
 H=-\omega_0^*J\omega_0,\qquad
 P=\gamma FE B^{-1},\\
 V&=X_0^*(P-\omega_0).
 \end{split}
\end{equation}
The hypotheses of the source theorem include \(\det B\ne0\).
The source already constructs exactly this \(\omega_0\) inside
\code{exists_isTriangularBasis}, then multiplies it by \(L\). The new proof stops
before that multiplication.

One has \(E^*D=\Id\), \(F^*J\omega_0=D=F^*JP\), and \(\omega_0\) lies in
the negative space \(yV_-\). Therefore \(H>0\), and the columns of
\(P-\omega_0\) are \(J\)-orthogonal to \(F\).
The existing
\code{gram_definiteCols_mul_eq} gives
\[
 V^*V=(P-\omega_0)^*J(P-\omega_0).
\]
Write \(K=\omega_0^*JP\). The span of \(\gamma F\) is negative semidefinite,
so \(P^*JP\leq0\). Expanding the preceding identity gives
\begin{equation}\label{eq:hermitian-lower}
 H+V^*V=P^*JP-K-K^*\leq-K-K^*.
\end{equation}
The determinant inequality for the Hermitian part of an accretive matrix then yields
\[
 2^{-(i+1)}\det(H+V^*V)\leq\abs{\det K}.
\]
For completeness, if \(\operatorname{Re}Q\geq M>0\), conjugating by \(M^{-1/2}\)
reduces to an accretive matrix whose eigenvalues have real part at least \(1\);
the product of their moduli is at least \(1\). Hence
\(\abs{\det Q}\geq\det M\). Apply this to \(Q=-K\) and
\(M=(H+V^*V)/2\).

The source's shared-vector determinant reduction, applied directly with \(D\),
gives
\[
 \nu=\code{volumeRatio}(\gamma,y)=\abs{\det B}\abs{\det K}.
\]
It uses \(E^*D=\Id\), the vanishing of \(D^*x\) on shared isotropic vectors,
and the existing
\code{norm_det_eq_norm_det_mul_inclMatrix}.
It does not require the general \code{IsNegativeBasis} structure.
Combining the last two displays gives the needed lower bound:
\begin{equation}\label{eq:raw-volume}
 \boxed{\quad
 \nu\ \geq\ \frac{\abs{\det B}}{2^{i+1}}
                   \det(H+V^*V).\quad}
\end{equation}

\subsection{The two geometric bounds replace every normalization factor}

Let \(r=i+1\), and suppose \(y\) satisfies the source bounds
\code{frameBox}\,\(\rho\) and \code{besselLength}\,\(y\leq\theta\).
Write the last column of \(V\) as \(v\), and \(h=H_{rr}\).
The needed estimates are exactly
\begin{equation}\label{eq:two-bounds}
 (\det H)^{-1}\leq\rho^{2r},\qquad
 (h+\norm v^2)^{-1}\leq\theta^{2(j-i)}.
\end{equation}

To prove the determinant estimate, let \(W=y^{-1}\omega_0\), so
\(W^*W=H\). For each subframe vector \(f_a\), its translate \(y^{-1}f_a\)
pairs with the columns of \(W\) as minus the \(a\)-th coordinate row.
Orthogonally projecting these translated subframe vectors onto the span of
\(W\) shows that \((H^{-1})_{aa}\leq\norm{y^{-1}f_a}^2\leq\rho^2\).
Hadamard's inequality for \(H^{-1}>0\) gives the first estimate.
An existing Hadamard lemma is available in
\src{Lib/General/LinearAlgebra/MinorProducts}{MinorProducts}.

For the last column, there are the same two geometric cases as in the source.
At a lower level, \(j=i\), the shared norm-one vector gives
\[
 h+\norm v^2=1.
\]
This follows directly by applying the existing shared-vector lemma to the final
column of \(P\), whose \(J\)-norm is \(-1\) and which is orthogonal to
\(P-\omega_0\). Expanding its Gram identity proves the displayed equality.

At the top level, \(j-i=1\). The final column \(w\) of \(\omega_0\) is orthogonal
to the isotropic part of the frame and has \(J\)-product \(1\) with the final
frame vector. Move \(y\) by the Bessel subgroup as in the existing
\code{norm_last_le} argument so that the relevant translated frame vector has
ordinary norm at most \(\theta\). Cauchy--Schwarz, using
\(\norm{y^{-1}w}^2=h\), gives
\[
 1\leq\theta\sqrt h.
\]
Since \(h>0\), this inequality also establishes \(\theta>0\).
Thus \(h\geq\theta^{-2}\), and adding \(\norm v^2\) gives the second estimate.
Neither argument constructs \(L\).

\subsection{The source constant is preserved exactly}

Under the source's affine relation, \(V=[U_{\rm base}+A^*,v]\), with \(H\) and
the final column fixed. Set \(\delta=\abs{\det B}\). By \eqref{eq:raw-volume},
\[
 \nu^{-(k+i)}
 \leq \left(\frac{2^{i+1}}{\delta}\right)^{k+i}
       \det\bigl(H+[U_{\rm base}+A^*,v]^*[U_{\rm base}+A^*,v]\bigr)^{-(k+i)}.
\]
Use the existing-Hua adapter and \eqref{eq:two-bounds}. The integral of this
majorant is at most
\begin{align*}
 \left(\frac{2^{i+1}}{\delta}\right)^{k+i}
  (\det H)^{-i}(h+\norm v^2)^{-k}
 &\leq
 \frac{2^{(i+1)(k+i)}\rho^{2i(i+1)}
                    \theta^{2k(j-i)}}{\delta^{k+i}}\\
 &=\frac{(2\rho^2)^{(i+1)i}
               (2^{i+1}\theta^{2(j-i)})^k}{\delta^{k+i}}.
\end{align*}
This is precisely the right-hand side of
\code{exists_majorant_volumeRatio_inv_pow}. Its hypotheses, measurable-majorant
conclusion, and constant can all remain unchanged. Consequently the downstream
\code{SchurOrbit} consumer need not change.
